% --- Archon physics formalization source begin ---
% archon:physics
% archon:covers PhyXMiniProblems/problem_phyx_mini_0115.lean
% archon:source-report reports/phyx_mini/problem_phyx_mini_0115.source.json

\chapter{Physics problem phyx\_mini\_0115}
\label{ch:PhyXMiniProblems_problem_phyx_mini_0115}

\paragraph{Problem source.}
\#\# Physical scenario

A 16.0-cm-long pencil is placed at a \$45.0\textasciicircum{}\textbackslash{}circ\$ angle, with its center 15.0 cm above the optic axis and 45.0 cm from a lens with a 20.0 cm focal length as shown in figure. (Note that the figure is not drawn to scale.) Assume that the diameter of the lens is large enough for the paraxial approximation to be valid.

\#\# Question

What is the length of the image?

\#\# Answer choices

- A: A: 17.1cm

- B: B: 18.2cm

- C: C: 19.3cm

- D: D: 16.8cm

\#\# Figure caption (auxiliary; use the image as primary evidence)

The image depicts a diagram involving a convex lens on a horizontal optical axis. Here's a detailed breakdown:

1. **Lens**: 
   - A convex lens is located on the right side, centered on the horizontal axis.

2. **Horizontal Axis**:
   - The optical axis runs horizontally through the center of the lens.

3. **Distances**:
   - A measurement of 45.0 cm is indicated on the horizontal axis from the leftmost vertical line (point A) to the center of the lens.

4. **Inclined Object**:
   - A line segment labeled as object \textbackslash{}( ABC \textbackslash{}) is positioned to the left of the lens, with three marked points A, C, and B.
   - The segment \textbackslash{}( BC \textbackslash{}) forms a 45.0° angle with the optical axis.
   - The vertical distance from point \textbackslash{}( A \textbackslash{}) to point \textbackslash{}( C \textbackslash{}) is labeled as 15.0 cm.

The diagram likely represents an optics problem illustrating the position and orientation of an object relative to a lens, and the associated distances and angles involved.

\#\# Dataset metadata

Geometrical Optics; Spatial Relation Reasoning, Physical Model Grounding Reasoning

\paragraph{Recorded answer/context.}
A: A: 17.1cm

\paragraph{Figure/image path.}
/root/proposal\_for\_physic/science-mango/phyx\_mini\_run/phyx\_data/test\_image/115.png

\paragraph{Formalization target.}
create a compiling Lean file with sorry bodies at `PhyXMiniProblems/problem\_phyx\_mini\_0115.lean`. The Lean declarations must preserve the physical quantities, dimensions or dimensional roles, figure labels, governing-law hypotheses, and final relation expressed by this problem.
Use Mathlib/Physlib names found through LeanExplore where available. If a physics API is missing, introduce faithful local abstractions rather than scalar placeholder aliases.

\begin{theorem}[Physics formalization target]
\label{thm:physics:phyx_mini_0115:target}
\lean{PhyXMiniProblems.ProblemPhyXMini0115.problem_phyx_mini_0115}
\uses{def:physics:phyx-mini-0115:phyxminiproblems-problemphyxmini0115-tiltedpencillenssetup, def:physics:phyx-mini-0115:phyxminiproblems-problemphyxmini0115-imagelengthincentimeters, def:physics:phyx-mini-0115:phyxminiproblems-problemphyxmini0115-matchespencillensfigure, def:physics:phyx-mini-0115:phyxminiproblems-problemphyxmini0115-hasadequateparaxialaperture, def:physics:phyx-mini-0115:phyxminiproblems-problemphyxmini0115-hasrealendpointimages, def:physics:phyx-mini-0115:phyxminiproblems-problemphyxmini0115-satisfiesparaxialthinlenslaws, def:physics:phyx-mini-0115:phyxminiproblems-problemphyxmini0115-roundstonearesttenthcentimeter, def:physics:phyx-mini-0115:phyxminiproblems-problemphyxmini0115-isclosestimagelengthanswer}
The assigned autoformalize agent should translate this physics problem into Lean declarations in the covered file, with theorem and lemma proof bodies written as `by sorry`.
\end{theorem}
\begin{proof}
This is an autoformalization task, not a proof task. Produce statements that can later be proved by the physics prover without weakening the physical model.
\end{proof}

% --- Archon declaration topology begin ---
\section{Declaration topology}
\begin{definition}[Length Quantity]
  \label{def:physics:phyx-mini-0115:phyxminiproblems-problemphyxmini0115-lengthquantity}
  \lean{PhyXMiniProblems.ProblemPhyXMini0115.LengthQuantity}
  A signed physical length represented coherently in every choice of units
\end{definition}
\begin{proof}
  This is the declared model definition of Length Quantity.
\end{proof}
\begin{definition}[centimeter Unit Choices]
  \label{def:physics:phyx-mini-0115:phyxminiproblems-problemphyxmini0115-centimeterunitchoices}
  \lean{PhyXMiniProblems.ProblemPhyXMini0115.centimeterUnitChoices}
  SI base units with centimeters selected as the length unit
\end{definition}
\begin{proof}
  This is the declared model definition of centimeter Unit Choices.
\end{proof}
\begin{definition}[length In Centimeters]
  \label{def:physics:phyx-mini-0115:phyxminiproblems-problemphyxmini0115-lengthincentimeters}
  \lean{PhyXMiniProblems.ProblemPhyXMini0115.lengthInCentimeters}
  \uses{def:physics:phyx-mini-0115:phyxminiproblems-problemphyxmini0115-lengthquantity, def:physics:phyx-mini-0115:phyxminiproblems-problemphyxmini0115-centimeterunitchoices}
  The signed scalar centimeter readout of a physical length
\end{definition}
\begin{proof}
  This is the declared model definition of length In Centimeters.
\end{proof}
\begin{definition}[Optical Plane Point]
  \label{def:physics:phyx-mini-0115:phyxminiproblems-problemphyxmini0115-opticalplanepoint}
  \lean{PhyXMiniProblems.ProblemPhyXMini0115.OpticalPlanePoint}
  \uses{def:physics:phyx-mini-0115:phyxminiproblems-problemphyxmini0115-lengthquantity}
  A point in the meridional plane of the lens, with dimensionful coordinates
\end{definition}
\begin{proof}
  This is the declared model definition of Optical Plane Point.
\end{proof}
\begin{definition}[axial Coordinate Cm]
  \label{def:physics:phyx-mini-0115:phyxminiproblems-problemphyxmini0115-axialcoordinatecm}
  \lean{PhyXMiniProblems.ProblemPhyXMini0115.axialCoordinateCm}
  \uses{def:physics:phyx-mini-0115:phyxminiproblems-problemphyxmini0115-lengthincentimeters, def:physics:phyx-mini-0115:phyxminiproblems-problemphyxmini0115-opticalplanepoint}
  Scalar axial-coordinate readout in centimeters
\end{definition}
\begin{proof}
  This is the declared model definition of axial Coordinate Cm.
\end{proof}
\begin{definition}[transverse Coordinate Cm]
  \label{def:physics:phyx-mini-0115:phyxminiproblems-problemphyxmini0115-transversecoordinatecm}
  \lean{PhyXMiniProblems.ProblemPhyXMini0115.transverseCoordinateCm}
  \uses{def:physics:phyx-mini-0115:phyxminiproblems-problemphyxmini0115-lengthincentimeters, def:physics:phyx-mini-0115:phyxminiproblems-problemphyxmini0115-opticalplanepoint}
  Scalar transverse-coordinate readout in centimeters
\end{definition}
\begin{proof}
  This is the declared model definition of transverse Coordinate Cm.
\end{proof}
\begin{definition}[plane Distance Cm]
  \label{def:physics:phyx-mini-0115:phyxminiproblems-problemphyxmini0115-planedistancecm}
  \lean{PhyXMiniProblems.ProblemPhyXMini0115.planeDistanceCm}
  \uses{def:physics:phyx-mini-0115:phyxminiproblems-problemphyxmini0115-opticalplanepoint, def:physics:phyx-mini-0115:phyxminiproblems-problemphyxmini0115-axialcoordinatecm, def:physics:phyx-mini-0115:phyxminiproblems-problemphyxmini0115-transversecoordinatecm}
  Euclidean separation of two meridional-plane points, read in centimeters
\end{definition}
\begin{proof}
  This is the declared model definition of plane Distance Cm.
\end{proof}
\begin{definition}[Pencil Point Label]
  \label{def:physics:phyx-mini-0115:phyxminiproblems-problemphyxmini0115-pencilpointlabel}
  \lean{PhyXMiniProblems.ProblemPhyXMini0115.PencilPointLabel}
  The three labels attached to the pencil in the primary figure
\end{definition}
\begin{proof}
  This is the declared model definition of Pencil Point Label.
\end{proof}
\begin{definition}[Pencil Endpoint]
  \label{def:physics:phyx-mini-0115:phyxminiproblems-problemphyxmini0115-pencilendpoint}
  \lean{PhyXMiniProblems.ProblemPhyXMini0115.PencilEndpoint}
  The two endpoints whose images determine the image length
\end{definition}
\begin{proof}
  This is the declared model definition of Pencil Endpoint.
\end{proof}
\begin{definition}[object Label]
  \label{def:physics:phyx-mini-0115:phyxminiproblems-problemphyxmini0115-pencilendpoint-objectlabel}
  \lean{PhyXMiniProblems.ProblemPhyXMini0115.PencilEndpoint.objectLabel}
  \uses{def:physics:phyx-mini-0115:phyxminiproblems-problemphyxmini0115-pencilpointlabel, def:physics:phyx-mini-0115:phyxminiproblems-problemphyxmini0115-pencilendpoint}
  Regard an endpoint label as its corresponding object-point label
\end{definition}
\begin{proof}
  This is the declared model definition of object Label.
\end{proof}
\begin{definition}[Thin Lens Kind]
  \label{def:physics:phyx-mini-0115:phyxminiproblems-problemphyxmini0115-thinlenskind}
  \lean{PhyXMiniProblems.ProblemPhyXMini0115.ThinLensKind}
  Optical classification of a thin lens
\end{definition}
\begin{proof}
  This is the declared model definition of Thin Lens Kind.
\end{proof}
\begin{definition}[Lens Profile]
  \label{def:physics:phyx-mini-0115:phyxminiproblems-problemphyxmini0115-lensprofile}
  \lean{PhyXMiniProblems.ProblemPhyXMini0115.LensProfile}
  Lens silhouette recorded by the primary figure
\end{definition}
\begin{proof}
  This is the declared model definition of Lens Profile.
\end{proof}
\begin{definition}[Optical Approximation]
  \label{def:physics:phyx-mini-0115:phyxminiproblems-problemphyxmini0115-opticalapproximation}
  \lean{PhyXMiniProblems.ProblemPhyXMini0115.OpticalApproximation}
  Approximation regime selected for geometrical imaging
\end{definition}
\begin{proof}
  This is the declared model definition of Optical Approximation.
\end{proof}
\begin{definition}[Thin Lens]
  \label{def:physics:phyx-mini-0115:phyxminiproblems-problemphyxmini0115-thinlens}
  \lean{PhyXMiniProblems.ProblemPhyXMini0115.ThinLens}
  \uses{def:physics:phyx-mini-0115:phyxminiproblems-problemphyxmini0115-lengthquantity, def:physics:phyx-mini-0115:phyxminiproblems-problemphyxmini0115-opticalplanepoint, def:physics:phyx-mini-0115:phyxminiproblems-problemphyxmini0115-thinlenskind, def:physics:phyx-mini-0115:phyxminiproblems-problemphyxmini0115-lensprofile}
  The pictured lens and its dimensional data. minimumDiameterForUnvignettedParaxialImage is the aperture threshold needed for the whole inclined pencil in the selected paraxial ray model. The problem does not provide a numerical aperture, only the statement that the actual diameter is at least this threshold
\end{definition}
\begin{proof}
  This is the declared model definition of Thin Lens.
\end{proof}
\begin{definition}[Tilted Pencil Lens Setup]
  \label{def:physics:phyx-mini-0115:phyxminiproblems-problemphyxmini0115-tiltedpencillenssetup}
  \lean{PhyXMiniProblems.ProblemPhyXMini0115.TiltedPencilLensSetup}
  \uses{def:physics:phyx-mini-0115:phyxminiproblems-problemphyxmini0115-lengthquantity, def:physics:phyx-mini-0115:phyxminiproblems-problemphyxmini0115-opticalplanepoint, def:physics:phyx-mini-0115:phyxminiproblems-problemphyxmini0115-pencilpointlabel, def:physics:phyx-mini-0115:phyxminiproblems-problemphyxmini0115-pencilendpoint, def:physics:phyx-mini-0115:phyxminiproblems-problemphyxmini0115-opticalapproximation, def:physics:phyx-mini-0115:phyxminiproblems-problemphyxmini0115-thinlens}
  The labeled pencil, its endpoint images, and the lens used to image it. No requested image length or answer choice is stored in this structure
\end{definition}
\begin{proof}
  This is the declared model definition of Tilted Pencil Lens Setup.
\end{proof}
\begin{definition}[object Distance Cm]
  \label{def:physics:phyx-mini-0115:phyxminiproblems-problemphyxmini0115-objectdistancecm}
  \lean{PhyXMiniProblems.ProblemPhyXMini0115.objectDistanceCm}
  \uses{def:physics:phyx-mini-0115:phyxminiproblems-problemphyxmini0115-axialcoordinatecm, def:physics:phyx-mini-0115:phyxminiproblems-problemphyxmini0115-pencilendpoint, def:physics:phyx-mini-0115:phyxminiproblems-problemphyxmini0115-tiltedpencillenssetup}
  Positive object distance from an object endpoint to the lens plane
\end{definition}
\begin{proof}
  This is the declared model definition of object Distance Cm.
\end{proof}
\begin{definition}[image Distance Cm]
  \label{def:physics:phyx-mini-0115:phyxminiproblems-problemphyxmini0115-imagedistancecm}
  \lean{PhyXMiniProblems.ProblemPhyXMini0115.imageDistanceCm}
  \uses{def:physics:phyx-mini-0115:phyxminiproblems-problemphyxmini0115-axialcoordinatecm, def:physics:phyx-mini-0115:phyxminiproblems-problemphyxmini0115-pencilendpoint, def:physics:phyx-mini-0115:phyxminiproblems-problemphyxmini0115-tiltedpencillenssetup}
  Positive real-image distance from the lens plane to an image endpoint
\end{definition}
\begin{proof}
  This is the declared model definition of image Distance Cm.
\end{proof}
\begin{definition}[object Height Cm]
  \label{def:physics:phyx-mini-0115:phyxminiproblems-problemphyxmini0115-objectheightcm}
  \lean{PhyXMiniProblems.ProblemPhyXMini0115.objectHeightCm}
  \uses{def:physics:phyx-mini-0115:phyxminiproblems-problemphyxmini0115-transversecoordinatecm, def:physics:phyx-mini-0115:phyxminiproblems-problemphyxmini0115-pencilendpoint, def:physics:phyx-mini-0115:phyxminiproblems-problemphyxmini0115-tiltedpencillenssetup}
  Signed object height relative to the optical axis through the lens center
\end{definition}
\begin{proof}
  This is the declared model definition of object Height Cm.
\end{proof}
\begin{definition}[image Height Cm]
  \label{def:physics:phyx-mini-0115:phyxminiproblems-problemphyxmini0115-imageheightcm}
  \lean{PhyXMiniProblems.ProblemPhyXMini0115.imageHeightCm}
  \uses{def:physics:phyx-mini-0115:phyxminiproblems-problemphyxmini0115-transversecoordinatecm, def:physics:phyx-mini-0115:phyxminiproblems-problemphyxmini0115-pencilendpoint, def:physics:phyx-mini-0115:phyxminiproblems-problemphyxmini0115-tiltedpencillenssetup}
  Signed image height relative to the optical axis through the lens center
\end{definition}
\begin{proof}
  This is the declared model definition of image Height Cm.
\end{proof}
\begin{definition}[image Length In Centimeters]
  \label{def:physics:phyx-mini-0115:phyxminiproblems-problemphyxmini0115-imagelengthincentimeters}
  \lean{PhyXMiniProblems.ProblemPhyXMini0115.imageLengthInCentimeters}
  \uses{def:physics:phyx-mini-0115:phyxminiproblems-problemphyxmini0115-planedistancecm, def:physics:phyx-mini-0115:phyxminiproblems-problemphyxmini0115-tiltedpencillenssetup}
  The requested physical image length, as a centimeter Euclidean readout
\end{definition}
\begin{proof}
  This is the declared model definition of image Length In Centimeters.
\end{proof}
\begin{definition}[Matches Pencil Lens Figure]
  \label{def:physics:phyx-mini-0115:phyxminiproblems-problemphyxmini0115-matchespencillensfigure}
  \lean{PhyXMiniProblems.ProblemPhyXMini0115.MatchesPencilLensFigure}
  \uses{def:physics:phyx-mini-0115:phyxminiproblems-problemphyxmini0115-lengthincentimeters, def:physics:phyx-mini-0115:phyxminiproblems-problemphyxmini0115-axialcoordinatecm, def:physics:phyx-mini-0115:phyxminiproblems-problemphyxmini0115-transversecoordinatecm, def:physics:phyx-mini-0115:phyxminiproblems-problemphyxmini0115-planedistancecm, def:physics:phyx-mini-0115:phyxminiproblems-problemphyxmini0115-tiltedpencillenssetup}
  Numerical and geometric readouts supplied by the statement and primary figure: the convex converging lens has focal length 20 cm; the 16 cm pencil has center C, which is 45 cm before the lens and 15 cm above its axis; and the directed segment A B is inclined at π/4 radians
\end{definition}
\begin{proof}
  This is the declared model definition of Matches Pencil Lens Figure.
\end{proof}
\begin{definition}[Has Adequate Paraxial Aperture]
  \label{def:physics:phyx-mini-0115:phyxminiproblems-problemphyxmini0115-hasadequateparaxialaperture}
  \lean{PhyXMiniProblems.ProblemPhyXMini0115.HasAdequateParaxialAperture}
  \uses{def:physics:phyx-mini-0115:phyxminiproblems-problemphyxmini0115-lengthincentimeters, def:physics:phyx-mini-0115:phyxminiproblems-problemphyxmini0115-tiltedpencillenssetup}
  The stated aperture assumption: the selected approximation is paraxial and the actual positive lens diameter meets the aperture threshold for an unvignetted image of both pencil endpoints
\end{definition}
\begin{proof}
  This is the declared model definition of Has Adequate Paraxial Aperture.
\end{proof}
\begin{definition}[Has Real Endpoint Images]
  \label{def:physics:phyx-mini-0115:phyxminiproblems-problemphyxmini0115-hasrealendpointimages}
  \lean{PhyXMiniProblems.ProblemPhyXMini0115.HasRealEndpointImages}
  \uses{def:physics:phyx-mini-0115:phyxminiproblems-problemphyxmini0115-lengthincentimeters, def:physics:phyx-mini-0115:phyxminiproblems-problemphyxmini0115-tiltedpencillenssetup, def:physics:phyx-mini-0115:phyxminiproblems-problemphyxmini0115-objectdistancecm, def:physics:phyx-mini-0115:phyxminiproblems-problemphyxmini0115-imagedistancecm}
  Both object endpoints lie beyond the focal plane and form real images
\end{definition}
\begin{proof}
  This is the declared model definition of Has Real Endpoint Images.
\end{proof}
\begin{definition}[Satisfies Paraxial Thin Lens Laws]
  \label{def:physics:phyx-mini-0115:phyxminiproblems-problemphyxmini0115-satisfiesparaxialthinlenslaws}
  \lean{PhyXMiniProblems.ProblemPhyXMini0115.SatisfiesParaxialThinLensLaws}
  \uses{def:physics:phyx-mini-0115:phyxminiproblems-problemphyxmini0115-lengthincentimeters, def:physics:phyx-mini-0115:phyxminiproblems-problemphyxmini0115-tiltedpencillenssetup, def:physics:phyx-mini-0115:phyxminiproblems-problemphyxmini0115-objectdistancecm, def:physics:phyx-mini-0115:phyxminiproblems-problemphyxmini0115-imagedistancecm, def:physics:phyx-mini-0115:phyxminiproblems-problemphyxmini0115-objectheightcm, def:physics:phyx-mini-0115:phyxminiproblems-problemphyxmini0115-imageheightcm}
  Endpointwise paraxial thin-lens laws. The first field is the homogeneous form f (p + q) = p q of the Gaussian lens equation. The second is the signed transverse-magnification law hᵢ p = -hₒ q. Neither field states the requested image length
\end{definition}
\begin{proof}
  This is the declared model definition of Satisfies Paraxial Thin Lens Laws.
\end{proof}
\begin{definition}[Answer Choice]
  \label{def:physics:phyx-mini-0115:phyxminiproblems-problemphyxmini0115-answerchoice}
  \lean{PhyXMiniProblems.ProblemPhyXMini0115.AnswerChoice}
  The four displayed multiple-choice labels
\end{definition}
\begin{proof}
  This is the declared model definition of Answer Choice.
\end{proof}
\begin{definition}[answer Image Length In Centimeters]
  \label{def:physics:phyx-mini-0115:phyxminiproblems-problemphyxmini0115-answerimagelengthincentimeters}
  \lean{PhyXMiniProblems.ProblemPhyXMini0115.answerImageLengthInCentimeters}
  \uses{def:physics:phyx-mini-0115:phyxminiproblems-problemphyxmini0115-answerchoice}
  Image-length readout printed beside each answer choice, in centimeters
\end{definition}
\begin{proof}
  This is the declared model definition of answer Image Length In Centimeters.
\end{proof}
\begin{definition}[Rounds To Nearest Tenth Centimeter]
  \label{def:physics:phyx-mini-0115:phyxminiproblems-problemphyxmini0115-roundstonearesttenthcentimeter}
  \lean{PhyXMiniProblems.ProblemPhyXMini0115.RoundsToNearestTenthCentimeter}
  \uses{def:physics:phyx-mini-0115:phyxminiproblems-problemphyxmini0115-answerchoice, def:physics:phyx-mini-0115:phyxminiproblems-problemphyxmini0115-answerimagelengthincentimeters}
  Agreement with a displayed image length rounded to the nearest tenth cm
\end{definition}
\begin{proof}
  This is the declared model definition of Rounds To Nearest Tenth Centimeter.
\end{proof}
\begin{definition}[Is Closest Image Length Answer]
  \label{def:physics:phyx-mini-0115:phyxminiproblems-problemphyxmini0115-isclosestimagelengthanswer}
  \lean{PhyXMiniProblems.ProblemPhyXMini0115.IsClosestImageLengthAnswer}
  \uses{def:physics:phyx-mini-0115:phyxminiproblems-problemphyxmini0115-answerchoice, def:physics:phyx-mini-0115:phyxminiproblems-problemphyxmini0115-answerimagelengthincentimeters}
  A selected answer is strictly closer to the exact image length than every rival
\end{definition}
\begin{proof}
  This is the declared model definition of Is Closest Image Length Answer.
\end{proof}
\begin{lemma}[object Endpoint Readouts exact]
  \label{lem:physics:phyx-mini-0115:phyxminiproblems-problemphyxmini0115-objectendpointreadouts-exact}
  \lean{PhyXMiniProblems.ProblemPhyXMini0115.objectEndpointReadouts_exact}
  \uses{def:physics:phyx-mini-0115:phyxminiproblems-problemphyxmini0115-tiltedpencillenssetup, def:physics:phyx-mini-0115:phyxminiproblems-problemphyxmini0115-objectdistancecm, def:physics:phyx-mini-0115:phyxminiproblems-problemphyxmini0115-objectheightcm, def:physics:phyx-mini-0115:phyxminiproblems-problemphyxmini0115-matchespencillensfigure}
  The figure geometry fixes the endpoint object distances and heights. In particular, half of the pencil has axial and transverse projections 4 √2 cm
\end{lemma}
\begin{proof}
  The conclusion follows from the governing relations and figure readouts named in the statement of object Endpoint Readouts exact.
\end{proof}
\begin{lemma}[image Endpoint Readouts exact]
  \label{lem:physics:phyx-mini-0115:phyxminiproblems-problemphyxmini0115-imageendpointreadouts-exact}
  \lean{PhyXMiniProblems.ProblemPhyXMini0115.imageEndpointReadouts_exact}
  \uses{def:physics:phyx-mini-0115:phyxminiproblems-problemphyxmini0115-tiltedpencillenssetup, def:physics:phyx-mini-0115:phyxminiproblems-problemphyxmini0115-imagedistancecm, def:physics:phyx-mini-0115:phyxminiproblems-problemphyxmini0115-imageheightcm, def:physics:phyx-mini-0115:phyxminiproblems-problemphyxmini0115-matchespencillensfigure, def:physics:phyx-mini-0115:phyxminiproblems-problemphyxmini0115-hasrealendpointimages, def:physics:phyx-mini-0115:phyxminiproblems-problemphyxmini0115-satisfiesparaxialthinlenslaws}
  Applying the Gaussian and magnification laws separately to A and B gives the exact axial distances and signed heights of A' and B'
\end{lemma}
\begin{proof}
  The conclusion follows from the governing relations and figure readouts named in the statement of image Endpoint Readouts exact.
\end{proof}
\begin{lemma}[image Length In Centimeters exact]
  \label{lem:physics:phyx-mini-0115:phyxminiproblems-problemphyxmini0115-imagelengthincentimeters-exact}
  \lean{PhyXMiniProblems.ProblemPhyXMini0115.imageLengthInCentimeters_exact}
  \uses{def:physics:phyx-mini-0115:phyxminiproblems-problemphyxmini0115-tiltedpencillenssetup, def:physics:phyx-mini-0115:phyxminiproblems-problemphyxmini0115-imagelengthincentimeters, def:physics:phyx-mini-0115:phyxminiproblems-problemphyxmini0115-matchespencillensfigure, def:physics:phyx-mini-0115:phyxminiproblems-problemphyxmini0115-hasrealendpointimages, def:physics:phyx-mini-0115:phyxminiproblems-problemphyxmini0115-satisfiesparaxialthinlenslaws}
  The endpoint image separation has exact length 3200 √10 / 593 cm
\end{lemma}
\begin{proof}
  The conclusion follows from the governing relations and figure readouts named in the statement of image Length In Centimeters exact.
\end{proof}
% --- Archon declaration topology end ---
% --- Archon physics formalization source end ---
